\chapter{Consideraciones Generales}

\subsection*{Dato clínico y enfoques complementarios}

Una vez realizado el par radiológico se pueden hacer enfoques complementarios, los cuales sirven para despejar dudas y aportar información diferente, el \textbf{dato clínico} determina cual es el enfoque complementario correspondiente en caso de ser necesario.\\

Los datos clínicos mas comunes son:
\begin{itemize}
    \item Traumatismo
    \item Fractura
    \item Luxación
    \item Dolor
    \item Derrame Articular
    \item Osteoporosis
    \item Gota
    \item Metástasis
    \item Tumores\\
\end{itemize}


\subsection*{Valor Diagnostico}
Se deben de cumplir ciertos requisitos para darle valor diagnostico a la radiografia, los cuales son:\\
\begin{itemize}
    \item Marca de la derecha
    \item Región comprendida
    \item Posición lograda
    \item Técnica adecuada\\
\end{itemize}

El valor diagnostico de la radiografía siempre esta relacionado al \textbf{dato clínico} \\

\begin{tcolorbox}[colback=red!10, colframe=red, title=Importante]
Una Rx cortada no tiene valor diagnostico ya que no cumple con la región comprendida, pero siempre \textbf{se debe de relacionar la Rx con su dato clínico}, para determinar si se le otorga o no valor.
\end{tcolorbox}

\subsection*{Parámetros de Exposición en Radiografía Convencional}

El \textbf{kilovoltaje pico (kVp)} determina la \textit{calidad} del haz de rayos X,
es decir, su poder de penetración. Un aumento del kVp produce fotones de mayor energía,
lo que amplía la escala de grises en la imagen resultante y, como consecuencia directa,
\textbf{disminuye el contraste}. A la inversa, valores bajos de kVp generan una escala
de grises corta, lo que se traduce en \textbf{alto contraste} —también denominado
``escala corta''—. El kVp influye además en la cantidad de radiación dispersa producida:
a mayor kVp, mayor dispersión, lo que deteriora la calidad de imagen si no se adoptan
medidas correctoras.

La \textbf{carga eléctrica (mAs)} es el parámetro que gobierna la \textit{cantidad} de
radiación, es decir, la dosis que recibe el paciente. Los mAs determinan el número de
electrones disponibles para la generación de rayos X: valores altos implican mayor
fluencia fotónica y, por tanto, mayor dosis; valores bajos la reducen. En la práctica,
kVp y mAs actúan de forma inversa sobre la dosis: el uso de \textbf{alto kVp} permite
reducir los mAs necesarios para obtener un nivel de ennegrecimiento adecuado,
disminuyendo así la dosis al paciente; mientras que el uso de \textbf{bajo kVp} obliga
a incrementar los mAs para compensar la menor penetración del haz, con el consiguiente
aumento de dosis.

La \textbf{rejilla antidifusora (Bucky)} se emplea cuando el espesor de la región
anatómica explorada genera una cantidad significativa de radiación dispersa. Al eliminar
los fotones dispersos antes de que alcancen el detector, la rejilla \textbf{mejora y
aumenta el contraste} de la imagen. Sin embargo, su uso implica un costo: la rejilla
absorbe también una fracción de la radiación primaria útil, por lo que exige un aumento
de la dosis de exposición para mantener el ennegrecimiento deseado. Además, la estructura
física de la rejilla introduce una leve reducción de la \textit{definición} espacial, ya
que sus láminas proyectan una trama sobre el receptor de imagen. Por esta razón, el uso
de Bucky solo está justificado cuando el espesor de la parte explorada realmente lo
amerita.\\

En la práctica clínica, una imagen \textbf{sobreexpuesta} —excesivamente ennegrecida—
puede corregirse \textit{reduciendo los mAs} o, si el contraste lo permite, \textit{aumentando el kVp}
para redistribuir la escala tonal. Una imagen \textbf{subexpuesta} —blanquecina y sin
detalle—, en cambio, se corrige principalmente \textit{aumentando los mAs}, ya que la cantidad
de radiación que llegó al detector fue insuficiente.

\subsection*{Marcación Radiológica}

La \textbf{marca radiológica} identifica el lado del paciente y debe colocarse
siguiendo criterios precisos según la proyección realizada. Como regla universal,
al montar la imagen en el negatoscopio siguiendo la convención estándar, la marca
queda invariablemente a la \textbf{izquierda del observador}, independientemente
de la proyección utilizada (salvo excepciones en algunos servicios).

En cuanto a la \textbf{altura}, la marca se coloca arriba cuando el paciente está
en bipedestación o a 45°, y abajo cuando está sentado o en decúbito supino.

En los \textbf{perfiles}, la posición física de la marca varía: el perfil derecho
se marca \textbf{anterior} y el perfil izquierdo se marca \textbf{posterior}. Cabe
destacar que en el \textbf{par radiológico} solo se marca el frente.

Algunas proyecciones presentan particularidades anatómicas que condicionan la
ubicación de la marca:

\begin{itemize}
    \item \textbf{Mano (PA):} la marca se coloca del lado del pulgar para la mano
    izquierda, y del lado del meñique para la mano derecha.
    \item \textbf{Antebrazo (AP):} la marca se ubica del lado del radio (externo)
    para el lado derecho, y del lado del cúbito (interno) para el lado izquierdo.
    \item \textbf{Pierna (AP):} la marca se coloca del lado del peroné (externo)
    para el lado derecho, y del lado de la tibia (interno) para el lado izquierdo.
    \item \textbf{Hombro (AP):} la marca se ubica en el lado externo para el derecho,
    y en el lado interno para el izquierdo.
    \item \textbf{Tórax (PA):} la marca queda del lado contrario a la silueta
    cardíaca.
\end{itemize}

\subsection*{Procedimientos Generales}

Previo a llamar al paciente, se debe alinear el tubo con el Bucky, luego con el
paciente y finalmente el rayo central (RC) en el punto de incidencia correspondiente.
Este procedimiento aplica tanto para el Bucky de pared como para la mesa.

\begin{enumerate}
    \item Llamar al paciente por nombre y apellido, verificando su identidad
    junto con la cédula de identidad.

    \item Hacerlo pasar, saludarlo y presentarse, dirigiéndose de forma cordial
    y profesional.

    \item Verificar que el pedido sea correcto: confirmar la región a estudiar y,
    en caso de que no especifique lateralidad, consultársela al paciente. Se puede
    indagar brevemente sobre el motivo de consulta.

    \item Medir el espesor de la región a explorar.

    \item Mientras el paciente se cambia, calcular la técnica radiográfica.

    \item Posicionar al paciente y darle las indicaciones previas a la exposición:
    no moverse, mantener la posición durante la apnea o respiración indicada, entre
    otras.

    \item Cerrar la bandeja del chasis.

    \item Cerrar las puertas del gabinete.

    \item Realizadas las exposiciones, indicar al paciente que puede relajarse y
    acomodarse según corresponda al estudio.
\end{enumerate}

\subsubsection*{Pregunta de Embarazo}

A toda paciente de sexo femenino en edad fértil ---comprendida aproximadamente
entre los 12 y 50 años--- se le debe consultar, por protocolo, si existe posibilidad
de embarazo. La pregunta se formula de forma discreta: \textit{``Disculpe, es por
protocolo: ¿hay alguna posibilidad de embarazo?''}

En caso de respuesta afirmativa, se le consultará si el médico solicitante fue
informado y si le explicó los riesgos asociados al estudio. Cuando la paciente
sea menor de edad y se encuentre acompañada por un familiar, se solicitará
amablemente que el acompañante retire por unos minutos antes de realizar la consulta.

\subsubsection*{Preparación del Paciente}

Se solicitará al paciente que retire todos los elementos que puedan generar
artefactos en la imagen:

\begin{itemize}
    \item Anillos y relojes. De no poder retirarlos, el paciente deberá lavarse
    con agua y jabón.
    \item Prendas con elementos metálicos: chapas, brillos, lentejuelas y
    estampados metálicos.
    \item Corpiños con aro, verificando siempre que el paciente haya cumplido
    con la indicación.
    \item Pantalones de jean, dado que la cremallera genera artefactos
    significativos en la imagen. Los pantalones deportivos o de tela pueden
    mantenerse puestos.
\end{itemize}

\subsubsection*{Consideraciones Prácticas}

\begin{itemize}
    \item \textbf{Paciente con dolor:} mantenerlo tranquilo, hablarle con calma e
    informarle que el estudio es rápido y necesario para el diagnóstico.

    \item \textbf{Paciente pediátrico acompañado por madre embarazada:} la madre
    debe retirarse. Ingresará otro acompañante, a quien se le colocará el chaleco
    plomado para permanecer durante la exposición. Es frecuente que los niños lloren
    durante el procedimiento; forma parte del manejo habitual.

    \item \textbf{Limpieza del chasis y la mesa radiográfica:} cuando se haya
    apoyado material potencialmente contaminado ---pie diabético, lesiones
    infectadas, entre otros--- se debe limpiar el chasis y la mesa con algodón
    embebido en alcohol. No se debe verter alcohol directamente sobre el chasis.
    Esta práctica debe aplicarse también durante el examen práctico.
    En lo posible limpiar siempre delante del paciente.

    \item \textbf{Uso de guantes:} si se utilizan guantes durante la atención al
    paciente, deben retirarse antes de manipular el tubo de rayos X.
\end{itemize}

\clearpage

\subsection*{Leyes de Proyección}

Las leyes de proyección describen el comportamiento geométrico del haz de rayos X
al interactuar con el objeto y el receptor de imagen, condicionando la fidelidad
y las limitaciones inherentes a toda imagen radiológica.

\begin{itemize}
    \item \textbf{Proyección cónica:} dado que los rayos X se originan en un foco
    puntual y se expanden en forma de abanico, la imagen resultante es siempre una
    proyección cónica. Como consecuencia, la imagen obtenida es invariablemente
    mayor que el objeto real, fenómeno conocido como \textit{magnificación}, aunque
    en condiciones estándar esta es mínima.

    \item \textbf{Superposición y sumación:} la imagen radiológica es la proyección
    de un objeto tridimensional sobre un plano bidimensional. Todas las estructuras
    atravesadas por el rayo central (RC) se aplanan en un único plano, generando
    superposición de estructuras y sumación de densidades. Un ejemplo representativo
    es la superposición de los cuerpos vertebrales con el contenido
    toracoabdominal y mediastínico.

    \item \textbf{Ley de las tangencias:} el contorno de una estructura solo se
    visualiza con nitidez cuando el RC incide de forma tangente a su superficie.
    Esta ley se cumple siempre. En la proyección de pelvis, la rotación interna de
    los fémures permite que el RC pase tangencial al cuello femoral, logrando
    visualizar sus bordes óseos con precisión. En la columna panorámica, se cumple
    en los bordes de los cuerpos vertebrales donde el RC logra pasar de forma
    tangente.

    \item \textbf{Ley de conservación de la forma:} para evitar la distorsión
    geométrica, el RC debe mantener una relación ortogonal ---a 90°--- con el
    receptor de imagen. Cuando esta condición se cumple, la forma del objeto se
    conserva de manera fiel en la imagen, con la única limitación de una
    magnificación mínima inherente a la proyección cónica.

    \item \textbf{Ángulo de incidencia:} esta ley se aplica cuando se requiere
    elongar una estructura o desplazar una superposición mediante la angulación
    deliberada del RC. En condiciones estándar, donde el RC incide de forma
    perpendicular al receptor, esta ley no se aplica y se prioriza la conservación
    de la forma anatómica.
\end{itemize}

\clearpage

\subsection*{Rúbrica de Evaluación Teórica}

\subsubsection*{Puntos Eliminatorios}

Los siguientes criterios son de carácter \textbf{eliminatorio}:

\begin{itemize}
    \item Reconocimiento del enfoque
    \item Posición lograda (incluyendo justificación y posibles mejoras)
    \item Determinar el valor diagnóstico de la imagen
\end{itemize}

\subsubsection*{Conocimientos de Mayor Complejidad}

\begin{table}[h]
\centering
\renewcommand{\arraystretch}{1.5}
\begin{tabular}{|p{3cm}|p{4.5cm}|p{4.5cm}|p{3.5cm}|}
\hline
\textbf{Criterio} & \textbf{25\%} & \textbf{10\%} & \textbf{0\%} \\
\hline
Reconocimiento del enfoque &
Reconoce el enfoque y si el DC es adecuado. Menciona par radiológico y enfoques complementarios. &
Reconoce el enfoque. No evalúa adecuación del DC ni menciona par radiológico o enfoques complementarios. &
No reconoce el enfoque ni evalúa ninguno de los aspectos anteriores. \\
\hline
Posición lograda &
Explica posición del paciente y RC con precisión. Determina si la posición está lograda mediante estructuras anatómicas proyectadas y justifica. Relaciona correctamente teoría y enfoque. &
Explica posición del paciente y RC de forma imprecisa. Determina posición lograda de forma vaga. Relaciona teoría y enfoque con imprecisión. &
Determina si la posición es adecuada sin relacionar con anatomía ni justificar, o no reconoce si está lograda. \\
\hline
Técnica adecuada &
Establece y justifica correctamente la técnica según el DC y la región. Reconoce si la técnica es adecuada en el enfoque entregado. &
Establece la técnica correctamente sin justificar y/o no reconoce con precisión si es adecuada en el enfoque. &
No establece correctamente la técnica ni reconoce si es adecuada. \\
\hline
\end{tabular}
\end{table}

\newpage
\subsubsection*{Conocimientos de Menor Complejidad}

\begin{table}[h]
\centering
\renewcommand{\arraystretch}{1.5}
\begin{tabular}{|p{3cm}|p{4.5cm}|p{4.5cm}|p{3.5cm}|}
\hline
\textbf{Criterio} & \textbf{12,5\%} & \textbf{5\%} & \textbf{0\%} \\
\hline
Región comprendida &
Conoce qué debe incluir la radiografía y lo señala. &
Conoce qué debe incluir pero no lo señala. &
No conoce qué debe incluir. \\
\hline
Marca de la derecha &
Reconoce si la radiografía está colocada correctamente según la marca y corrige si es necesario (derecha-izquierda, arriba-abajo). Reconoce lateralidad en miembro o perfil si corresponde. &
Falta precisión en el reconocimiento de la colocación según la marca. No reconoce lateralidad si corresponde. &
No reconoce la ubicación de la radiografía según la marca. \\
\hline
\end{tabular}
\end{table}

\clearpage

\subsection*{Examen: Enfoques y DC Frecuentes}

Al ingresar al examen, las tres imágenes ya se encuentran colocadas y el DC
es comunicado de forma oral por el evaluador. Se recomienda seguir estrictamente
los puntos indicados en la rúbrica para evitar preguntas complementarias.

A continuación se listan los enfoques y diagnósticos clínicos (DC) relevados
como frecuentes en instancias previas (en negrita los que mas se repiten):

\begin{multicols}{2}
\begin{itemize}
    \item \textbf{Frente de tórax con catéter venoso central}
    \item Hueso nasal perfil
    \item \textbf{Vértices pulmonares (DC: tuberculosis)}
    \item Perfil de sacro-cóccix
    \item Maxilar inferior desenfilado
    \item \textbf{Perfil de tórax}
    \item Perfil de silla turca
    \item Frente de columna cervical (mal posicionado, DC: traumatismo)
    \item Outlet
    \item \textbf{Oblicua de vejiga}
    \item Caldwell modificado
    \item Frente de lumbosacra (tornillos)
    \item Columna lumbosacra perfil
    \item Frente de cráneo (mal posicionado)
    \item Funcional de lumbar en flexión (hiperflexión)
    \item \textbf{Abdomen de pie (DC: oclusión intestinal)}
    \item Waters (DC: traumatismo)
    \item Pelvis AP común (prótesis derecha)
    \item Cigomático unilateral
    \item Hirtz modificado para cigomático
    \item Frente de sacro
    \item Caldwell a 15°
    \item \textbf{Frente de abdomen}
    \item \textbf{Pie (DC: traumatismo)}
    \item Perfil de columna cervical
    \item Waters-Mahoney
    \item \textbf{Abdomen en decúbito}
    \item Panorámica
    \item Schüller
    \item \textbf{Omalgia (ambas rotaciones)}
    \item Desfiladero
    \item Alar y obturatriz
\end{itemize}
\end{multicols}

\begin{tcolorbox}[colback=red!10, colframe=red, title=Importante]
Si te preguntan si se puede mejorar, siempre di que si, ya que siempre algo en la Rx se puede mejorar
\end{tcolorbox}